\documentclass[tikz,border=4pt]{standalone}
\usepackage{ctex}
\usepackage{amsmath}
\usetikzlibrary{arrows.meta}

% p8-02e: 一次函数例题——判断 y=-2x+3 经过哪些象限（k<0, b>0）
\begin{document}
\begin{tikzpicture}[scale=0.9, thick]
  % 坐标轴
  \draw[->] (-2.0,0) -- (2.5,0) node[right] {$x$};
  \draw[->] (0,-2.0) -- (0,4.0) node[above] {$y$};
  \node[below left] at (0,0) {$O$};

  % 刻度
  \foreach \x in {-1,1,2} {
    \draw (\x, 2pt) -- (\x, -2pt) node[below] {\small$\x$};
  }
  \foreach \y in {-1,1,2,3} {
    \draw (2pt,\y) -- (-2pt,\y) node[left] {\small$\y$};
  }

  % 直线 y = -2x + 3
  \draw[blue, thick] (-0.5, 4.0) -- (2.5, -2.0);
  \node[right, blue] at (2.1, -1.5) {$y=-2x+3$};

  % 纵截距和横截距
  \fill[blue] (0,3) circle (2.5pt);
  \node[left, blue] at (0,3) {$(0,3)$};

  \fill[blue] (1.5,0) circle (2.5pt);
  \node[above right, blue] at (1.5,0) {$\left(\frac{3}{2},0\right)$};

  % 象限标注
  \node at (1.5,1.5) {\small I};
  \node at (-1.2,1.5) {\small II};
  \node at (-1.2,-1.2) {\small III};
  \node at (1.5,-1.2) {\small IV};

  % 标注 k<0, b>0
  \node[blue, font=\small] at (-1.5, 3.5) {$k=-2<0$};
  \node[blue, font=\small] at (-1.5, 2.8) {$b=3>0$};
  \node[blue, font=\small] at (-1.5, 2.1) {经过I,II,IV象限};
\end{tikzpicture}
\end{document}
